\section{Embedded limits and the punctual genus correction}
\label{sec:embedded-limits-v169}
The parameter fibre of $T_a$ is a reduced chain; its universal curves
need not be reduced or Cohen--Macaulay. These are different assertions
about different schemes. We now determine both the one-dimensional
parts and the embedded zero-dimensional structure of every curve over
that chain.

\subsection{All formal arcs in the marked family}
Let $b=\beta\tau^p$, $c=\gamma\tau^q$, with $p,q>0$ and
$\beta,\gamma\in k[[\tau]]^*$. The units may be arbitrary formal
series. The following ideals are written on $s=1$ with homogeneous
target coordinates. At source infinity they are glued to $G=R=0$;
equivalently, take the ideal sheaves specified by
Proposition~\ref{prop:curve-charts-v169}.

\begin{theorem}[The complete arc table]\label{thm:arc-table-v169}
There are $a$ walls $p=jq$, $1\leq j\leq a$, and $a+1$ open chambers.
The special embedded ideals are as follows:
\begin{align}
 \mathcal O_0\ (p<q):\quad& (R,x^aG),\label{eq:limit-O0-v169}\\
 \mathcal W_1\ (p=q):\quad&
       (R-\kappa x^{a-1}G,x^aG),\quad
       \kappa=\overline\gamma/\overline\beta,\label{eq:limit-W1-v169}\\
 \mathcal O_i\ (iq<p<(i+1)q):\quad&
       (xR,R^{i+1},x^{a-j}F^{j-1}G\ (1\leq j\leq i)),
       \quad1\leq i<a,\label{eq:limit-Oi-v169}\\
 \mathcal W_j\ (p=jq):\quad&
       (xR,x^{a-r}F^{r-1}G\ (1\leq r<j),
                   R^j-\lambda x^{a-j}F^{j-1}G),
       \quad2\leq j\leq a,\label{eq:limit-Wj-v169}\\
 \mathcal O_a\ (p>aq):\quad&
       (xR,x^{a-j}F^{j-1}G\ (1\leq j\leq a)).
       \label{eq:limit-Oa-v169}
\end{align}
In \eqref{eq:limit-Wj-v169}, $\lambda=\overline\gamma^j/\overline\beta$.
Each wall is a one-dimensional family with nonzero residue parameter;
each open chamber gives one fixed embedded curve. The closure of the
$j$th wall family is $E_j\simeq\Pj^1$, with endpoints
$\mathcal O_{j-1},\mathcal O_j$. These are all components, points, and
adjacencies of the parameter fibre of $T_a$. There is only one
normalization branch at every point of this smooth parameter surface.
The fan is minimal for these fixed-target embedded limits.
\end{theorem}
\begin{proof}
In the first chart $k=c/b$ has nonnegative valuation exactly when
$p\leq q$. Its residue is zero in the first chamber and is $\kappa$
on the first wall. In $T_i$ the two coordinates have valuations
$p-iq$ and $(i+1)q-p$. Setting their positive-valuation coordinates
to zero in \eqref{eq:curve-middle-v169} gives
\eqref{eq:limit-Oi-v169}. Setting $e=0$ and retaining
$h=\lambda\ne0$ gives the next wall; the lower wall is its expression
in the adjacent chart. In the last chart $z=b/c^a$ has nonnegative
valuation exactly when $p\geq aq$. Its nonzero residue gives
$\mathcal W_a$ with $\lambda=\overline z^{-1}$; $z=0$ gives
\eqref{eq:limit-Oa-v169}. Properness and the chart cover exhaust all
arcs. If $c$ is identically zero, the first chart gives
\eqref{eq:limit-O0-v169} as well.

Every nonzero wall residue is attained by an arc with $p=jq$.
The transition formulas identify its two endpoints as stated. The
Newton facets in Proposition~\ref{prop:rees-v169} are all present;
equivalently, the residue term in \eqref{eq:limit-W1-v169} or
\eqref{eq:limit-Wj-v169} changes the ideal sheaf as the residue varies.
A wall cannot therefore be removed while retaining the embedded limit
map. This minimal fan belongs to the marked two-parameter problem,
not to an arbitrary coefficient presentation on $B_a$.
\end{proof}

\begin{theorem}[A sharp, degree-independent jet bound on the family]
\label{thm:sharp-jet-v169}
Let a formal arc in $S_a$ be centred at the origin and satisfy
$\operatorname{ord}(b)=p<\infty$. Its closed embedded Hilbert limit is
determined by the coefficient arc modulo $\tau^{p+1}$. This bound
is independent of $a$ and of the Hilbert embedding degree. For every
$p\geq1$ it is optimal as a uniform sufficient bound over arcs with
that value of $p$.
\end{theorem}
\begin{proof}
If $\operatorname{ord}(c)>p$, including $c=0$, any congruent arc is
still in $\mathcal O_0$. Otherwise congruence preserves $p,q$ and the
leading coefficients of both $b,c$. It preserves every wall residue
in Theorem~\ref{thm:arc-table-v169}, which proves sufficiency.
For necessity take $c=\tau^p$ and $b=\beta\tau^p$ for two different
nonzero constants $\beta$. The arcs agree modulo $\tau^p$, but the
first-wall ideals differ: on $G=1$, the nonzero class $x^{a-1}$ has
$R=\beta^{-1}x^{a-1}$ and $x^a=0$. Thus that jet does not determine
the fixed-target limit.
\end{proof}
Here $p$ is the intersection order with the marked base-point divisor
$b=0$. It is not the raw Smith length. In degree $m$ the latter is
\begin{equation}\label{eq:raw-order-v169}
 \mu=Cp+\sum_{j=1}^aE_j\min(p,jq),
\end{equation}
whereas the primitive degree-$m$ ideal has order given by the sum
without $Cp$, and the unweighted graph ideal has order
$\sum_{j=1}^a\min(p,jq)$. These are three distinct quantities, none
of which should be substituted for the optimal family bound $p+1$.
For $a=3,m=7$, arcs $(b,c)=(\tau,\tau^2)$ and
$(\tau^2,\tau^3)$ have the same special Hilbert curve but raw Smith
lengths $84$ and $168$. Ramifying an arc multiplies all ideal orders,
while leaving its embedded closed limit unchanged.

\subsection{Local rings and associated points}
Let $L_0$ denote the main component $G=R=0$. Let $L$ denote the tail
$x=R=0$, and put
\[
 P=(x=0,[F:G:R]=[0:1:0]).
\]
All local ideals below are ordinary ideals in the indicated affine
rings, before passage to sheaves. Write $H^0_P(\OO_C)$ for the
largest submodule supported at the closed point $P$, not for the
whole nilradical of a curve with a generically thick tail.

\begin{theorem}[Pure tails and punctual excess]\label{thm:punctual-v169}
For $\mathcal O_0$ and $\mathcal W_1$ the curve is Cohen--Macaulay,
with components $L_0,L$ of generic multiplicities $1,a$.
For $1\leq i<a$, the curve $\mathcal O_i$ has the same two minimal
components and multiplicities. On $F=1$ its ideal is
\begin{equation}\label{eq:Oi-F-v169}
 (xR,x^{a-i}G,R^{i+1})
    =(G,R)\cap(xR,x^{a-i},R^{i+1}).
\end{equation}
On $G=1$ it is
\begin{equation}\label{eq:Oi-G-v169}
 (xR,R^{i+1},x^{a-i}(x,F)^{i-1}).
\end{equation}
There is precisely one embedded associated point when $i\geq2$, namely
$P$, and none when $i=1$. Its punctual module is
\begin{equation}\label{eq:Oi-torsion-v169}
 H^0_P(\OO_{\mathcal O_i})
   \simeq x^{a-i}k[x,F]\big/x^{a-i}(x,F)^{i-1},\qquad R\cdot H^0_P=0,
 \qquad\length H^0_P=\binom i2.
\end{equation}

For $2\leq j\leq a$ put $d=a-j$. The wall curve $\mathcal W_j$ has
on $F=1$ the complete-intersection ideal
\begin{equation}\label{eq:W-F-v169}
 (xR,R^j-\lambda x^dG).
\end{equation}
On $G=1$ its ideal is
\begin{equation}\label{eq:W-G-v169}
 (xR,x^{d+1}(x,F)^{j-2},R^j-\lambda x^dF^{j-1}).
\end{equation}
Its pure quotient on this chart is obtained by adding $x^{d+1}$.
The kernel is
\begin{equation}\label{eq:W-torsion-v169}
 H^0_P(\OO_{\mathcal W_j})
  \simeq x^{d+1}k[x,F]\big/x^{d+1}(x,F)^{j-2},\qquad R\cdot H^0_P=0,
 \qquad\length H^0_P=\binom{j-1}{2}.
\end{equation}
For $j<a$ the minimal components are again $L_0,L$ with multiplicities
$1,a$. For $j=a$ the tail is the reduced plane curve
$R^a=\lambda F^{a-1}G$. In both cases the only possible embedded point
is $P$, and it occurs exactly when $j\geq3$.
\end{theorem}
\begin{proof}
Eliminating $R$ in the first two cases leaves a hypersurface, and gives
the multiplicities directly. For $\mathcal O_i$, localization at $F$
reduces all the monomials in \eqref{eq:limit-Oi-v169} to
\eqref{eq:Oi-F-v169}. The intersection identity follows by monomial
divisibility. The second component is primary to $(x,R)$ and is free
of rank $a$ over $k[G]$, with transverse basis
\[
 1,x,\ldots,x^{a-i-1},R,\ldots,R^i.
\]
The entire ring in \eqref{eq:Oi-F-v169} has a length-two
Hilbert--Burch resolution; alternatively its three monomial generators
have the two syzygies coming from $xR^{i+1}$ and $x^{a-i}GR$.
It is one-dimensional Cohen--Macaulay and has no embedded point.

On $G=1$ the ideal is \eqref{eq:Oi-G-v169}. Adding $x^{a-i}$ gives
the pure primary ideal $(xR,R^{i+1},x^{a-i})$. The kernel is generated
by the image of $x^{a-i}$, with annihilator
$(R,(x,F)^{i-1})$ on that generator. This proves
\eqref{eq:Oi-torsion-v169}. For $i\geq2$ an explicit primary
decomposition is
\[
 (xR,R^{i+1},x^{a-i})
 \ \cap\ (xR,R^{i+1},(x,F)^{a-1}).
\]
The second factor is primary to $(x,F,R)$. The displayed length is
the number of monomials in two variables of total degree below $i-1$.

For a wall, the equations reduce to \eqref{eq:W-F-v169} and
\eqref{eq:W-G-v169}. The former is a regular sequence of length two.
In the latter, after adjoining $x^{d+1}$, there is a free $k[F]$-basis
\[
 1,x,\ldots,x^d,R,\ldots,R^{j-1};
\]
the relation $R^j=\lambda F^{j-1}x^d$ and $xR=0$ give its
multiplication. Hence this pure ring is Cohen--Macaulay of rank $a$
over $k[F]$. The kernel of the pure quotient is generated by
$x^{d+1}$ and has exactly the annihilator described in
\eqref{eq:W-torsion-v169}. No extra relation arises from multiplying
$R^j-\lambda x^dF^{j-1}$ by $x$: the resulting
$x^{d+1}F^{j-1}$ already lies in $x^{d+1}(x,F)^{j-2}$.
A monomial reduction in $x,F$, followed by reduction in $R$, proves
that these are all the relations. This also computes the length.
For $d>0$ the pure ideal is primary to $(x,R)$; for $d=0$ it is the
reduced cuspidal tail ideal $(x,R^a-\lambda F^{a-1})$.

There are no missing associated points. Away from $P$ the $F$ chart
or its overlap has just been treated. On $R=1$ the nonterminal cases
have empty special fibre, while on the last wall $F,G$ are units and
\eqref{eq:W-F-v169} applies. Source infinity is
\eqref{eq:source-infinity-v169}. These charts cover the curve.
\end{proof}

\begin{theorem}[The cuspidal wall and its nilradical]\label{thm:genus-correction-v169}
On the last wall, $\mathcal W_a$ consists of its reduced curve
$L_0\cup Q_a$ together with a punctual nilradical at $P$, where
\[
 Q_a:\ x=0,\quad R^a=\lambda F^{a-1}G.
\]
The normalization of $Q_a$ is $\Pj^1$, with map
\begin{equation}\label{eq:cusp-normalization-v169}
 [u:v]\longmapsto[u^a:\lambda^{-1}v^a:u^{a-1}v].
\end{equation}
The main component meets $Q_a$ transversely at its smooth point
$[1:0:0]$. Put $g_a=(a-1)(a-2)/2$. The nilradical $\mathcal N$ has
length $g_a$ and, for $a\geq3$, nilpotence order $a-1$. More precisely,
on $G=1$ it is the ideal generated by $x$ with
\begin{equation}\label{eq:genus-module-v169}
 \mathcal N\simeq xk[x,F]/x(x,F)^{a-2},\quad R\mathcal N=0,
 \qquad
 \length(\mathcal N^r)=\binom{a-r}{2}\quad(1\leq r\leq a-2).
\end{equation}
For $a=2$ it is zero. The Hilbert polynomial of the reduced curve is
$(a+1)l+1-g_a$; this exact nilradical restores $P_a(l)=(a+1)l+1$.
Thus the punctual part is forced by arithmetic genus, not removable
excess in a flat Hilbert limit.
\end{theorem}
\begin{proof}
Set $j=a,d=0$ in \eqref{eq:W-G-v169}. The ring is
\[
 k[x,F,R]/(xR,x(x,F)^{a-2},R^a-\lambda F^{a-1}).
\]
Its reduction is the domain $k[F,R]/(R^a-\lambda F^{a-1})$; the
nilradical is exactly $(x)$. The formula for its module follows from
Theorem~\ref{thm:punctual-v169}. A basis of its $r$th power consists
of $x^uF^v$ with $u\geq r$ and $u+v\leq a-2$. Counting these
monomials gives \eqref{eq:genus-module-v169} and proves the
nilpotence statement, including nonvanishing of $x^{a-2}$.
The parametrization \eqref{eq:cusp-normalization-v169} is birational;
on the cusp chart its local semigroup is generated by $a-1,a$.
Its missing nonnegative exponents number $g_a$. The only singularity
of the plane tail is that cusp, while at the attaching point $F=1$
the full reduced local ring is $k[x,R]/(xR)$ after eliminating $G$.
The normalization of the whole reduction is consequently two copies
of $\Pj^1$. A plane curve of degree $a$ has arithmetic genus $g_a$;
attaching the main line at one point gives the stated reduced Hilbert
polynomial. The exact sequence
$0\to\mathcal N\to\OO_{\mathcal W_a}\to\OO_{(\mathcal W_a)_{\rm red}}\to0$
then proves the correction formula.
\end{proof}

\begin{proposition}[The terminal chamber]\label{prop:terminal-v169}
The curve $\mathcal O_a$ has three minimal components: $L_0$ with
multiplicity one, the tail $x=G=0$ with multiplicity one, and the
tail $x=F=0$ with multiplicity $a-1$. Their reduced supports form a
chain of three lines. On $F=1$ the ideal is $(G,xR)$; on $G=1$ it is
\[
 (xR,(x,F)^{a-1}).
\]
The latter has pure quotient $(x,F^{a-1})$ and punctual kernel
$xk[x,F]/x(x,F)^{a-2}$ of length $g_a$. On $R=1$ the ideal is
$(x,F^{a-1}G)$, a hypersurface with no embedded point. Thus $P$ is
the only embedded associated point when $a\geq3$, and there is none
when $a=2$. The full nilradical is not just this punctual module:
the $F=0$ tail is generically nonreduced for $a>2$.
\end{proposition}
\begin{proof}
Localize \eqref{eq:limit-Oa-v169} in each target coordinate. On $G=1$
the monomials $x^{a-j}F^{j-1}$ generate $(x,F)^{a-1}$.
Adding $x$ yields the pure ring, free of rank $a-1$ over $k[R]$;
the kernel has the stated two-variable basis and is killed by $R$.
On $R=1$, $xR$ eliminates $x$ and leaves $F^{a-1}G$.
These ordinary affine calculations give the minimal multiplicities
and all associated points. The normalization of the reduction is the
disjoint union of its three lines. The pure plane tail
$F^{a-1}G=0$ has arithmetic genus $g_a$, so the same length correction
also recovers $P_a$ in this terminal chamber.
\end{proof}

\subsection{The first higher-contact instance}
For $a=3$ the Gotzmann degree is $7$, and
\begin{equation}\label{eq:cubic-factor-v169}
 I_{29}(M_{3,7})\OO_{S_3}
       =b^{56}(b,c)^7(b,c^2)^6(b,c^3)^{15}.
\end{equation}
The parameter fibre has three components, with self-intersections
$-2,-2,-1$. Every one is one-dimensional, their successive
intersections are single reduced points, and there are no other
intersections or normalization branches. The seven embedded curve
families are collected here; each ideal is on $s=1$.
\begin{center}
\small
\begin{tabular}{c|l|c}
region & ideal & $\length H^0_P$\\ \hline
$p<q$ & $(R,x^3G)$ & $0$\\
$p=q$ & $(R-\kappa x^2G,x^3G)$ & $0$\\
$q<p<2q$ & $(xR,x^2G,R^2)$ & $0$\\
$p=2q$ & $(xR,x^2G,R^2-\lambda xFG)$ & $0$\\
$2q<p<3q$ & $(xR,x^2G,xFG,R^3)$ & $1$\\
$p=3q$ & $(xR,x^2G,xFG,R^3-\lambda F^2G)$ & $1$\\
$p>3q$ & $(xR,x^2G,xFG,F^2G)$ & $1$
\end{tabular}
\end{center}
At the last wall this length-one module is square-zero and is supported
at the cusp of the cubic tail. At the preceding open chamber it is
an embedded point on a generically triple line. These are different
local algebras with the same required Euler-characteristic correction.
For $a\geq4$ the last-wall nilradical is no longer square-zero, as
\eqref{eq:genus-module-v169} shows. This is why the first-contact
extension algebra alone cannot describe higher contacts.

The entire $S_3$ parameter fibre has been classified above, including
its universal curves. The full $B_3$ parameter fibre has more
coefficient directions, and no equality with this three-component
chain is asserted. The distinction remains important when the
marked family is inserted into a local Smith normal form of a pencil.
